\chapter{RiemannPartitionUp NameCert construction}
\label{ch:concrete-instances-riemann-partition-namecert}
\origin{human}

Following Riemann and Darboux, the $\RiemannPartitionUp$ packet records a finite subdivision of a displayed $\RealUp$ interval as NameCert data before tags, sums, or integrability claims are available. It sits after the dyadic partition route of \autoref{ch:concrete-instances-dyadic-partition-namecert}, the equal-mesh route of \autoref{ch:concrete-instances-uniform-partition-namecert}, and the finite mesh-refinement surface of \autoref{ch:concrete-instances-finite-partition-mesh-namecert}. The packet gives L10 interval consumers a standard finite-subdivision interface without importing host finite lists, quotient partition equality, selected mesh limits, or ambient $\RealUp$ completeness.

\begin{definition}[RiemannPartition carrier]
\label{def:riemann-partition-carrier}
A $\RiemannPartitionUp$ carrier is a finite $\BHist$ packet
$$
P=(I,D,U,M,R,H,C,Q,N)
$$
with the following displayed rows:
\begin{enumerate}
\item $I$ is the $\RealUp$ interval source row, carrying the named left and right endpoint seals through which the subdivision is read.
\item $D$ is the $\DyadicPartitionUp$ row, listing the dyadic breakpoint cells that refine the displayed interval.
\item $U$ is the $\UniformPartitionUp$ row, recording any equal-mesh window used by the same subdivision request.
\item $M$ is the $\FinitePartitionMeshUp$ row, exposing the mesh bound, adjacency, endpoint coverage, and common-refinement witness for the displayed cells.
\item $R$ is the downstream readback row for Darboux, Riemann-sum, regulated-function, or interval-estimate consumers that inspect the subdivision only after $I,D,U,M$ have been displayed.
\item $H$ records componentwise $\hsame$ transport for interval endpoints, dyadic cells, uniform mesh data, finite mesh witnesses, downstream readback, replay, provenance, and local naming rows.
\item $C$ records displayed $\Cont$ replay from an interval-subdivision request through $I,D,U,M,R,H,Q,N$.
\item $Q$ is the $\Pkg$ provenance over $\RiemannPartitionUp$, $\RealUp$, $\DyadicPartitionUp$, $\UniformPartitionUp$, $\FinitePartitionMeshUp$, $\BHist$, $\hsame$, $\Cont$, and $\NameCert$.
\item $N$ is the local $\NameCert_{\RiemannPartitionUp}$ row.
\end{enumerate}
The classifier compares accepted packets only through $I,D,U,M,R,H,C,Q,N$. It has no coordinate for tagged Riemann sums, Darboux equality, host finite-list equality, quotient partitions, selected mesh limits, countable choice, or ambient $\RealUp$ completeness.
\end{definition}

\begin{theorem}[RiemannPartition mesh refinement]
\label{thm:riemann-partition-mesh-refinement}
Every accepted $\RiemannPartitionUp$ carrier $P=(I,D,U,M,R,H,C,Q,N)$ refines its displayed $\RealUp$ interval through the existing $\DyadicPartitionUp$, $\UniformPartitionUp$, and $\FinitePartitionMeshUp$ carrier rows before any subdivision consumer can read the packet. The refinement route exposes dyadic cells through $D$, equal-mesh windows through $U$ when present, and mesh-bound and common-refinement witnesses through $M$; it exports no Riemann integrability theorem, Darboux equality, host finite-list presentation, quotient partition equality, or selected limiting partition.
\end{theorem}
\begin{proof}
Project the accepted carrier of \autoref{def:riemann-partition-carrier}. The interval source $I$ is the only endpoint row, while $D,U,M$ are the only subdivision and mesh rows available to a consumer.

Read the dyadic side through \autoref{ch:concrete-instances-dyadic-partition-namecert}. The row $D$ supplies the finite dyadic cells and their ordered interval coverage, so a consumer cannot replace the subdivision by a host list of breakpoints.

Next read the equal-mesh and mesh-refinement rows through \autoref{ch:concrete-instances-uniform-partition-namecert} and \autoref{ch:concrete-instances-finite-partition-mesh-namecert}. These rows expose uniform windows, mesh bounds, adjacency, endpoint coverage, and common-refinement data before $R$ is public.

Finally replay the structural rows $H,C,Q,N$. Transport, continuation replay, provenance, and local naming preserve the same finite subdivision packet, so the accepted route remains a NameCert interface for interval work rather than an integrability theorem, Darboux equality, quotient partition principle, selected mesh-limit construction, or host finite-list assertion.
\end{proof}

\falsifiablePrediction{A $\RiemannPartitionUp$ consumer that reads a subdivision without exposing the $\RealUp$ interval row, dyadic partition row, optional uniform-mesh row, finite mesh-refinement row, transport, replay, provenance, and local $\NameCert$ row refutes the carrier surface.}

\independenceWitness{dyadic-partition, uniform-partition, finite-partition-mesh: $\RiemannPartitionUp$ is not bijective with any listed sibling because it joins the interval source, dyadic cells, optional uniform window, and common mesh-refinement handoff as one subdivision packet for downstream interval consumers.}

\closureat{RiemannPartitionUp}{seedStr}

\begin{closurestatus}{\RiemannPartitionUp}
  \constructivestory{A $\RiemannPartitionUp$ packet is generated from a RealUp interval row, DyadicPartition cells, optional UniformPartition mesh windows, FinitePartitionMesh refinement witnesses, downstream readback, componentwise hsame transport, displayed Cont replay, Pkg provenance, and a local NameCert row.}
  \theoryclosure{\seedClosure}
  \scopeclosed{\autoref{def:riemann-partition-carrier} and \autoref{thm:riemann-partition-mesh-refinement} seed the finite Real interval subdivision interface through DyadicPartitionUp, UniformPartitionUp, and FinitePartitionMeshUp rows.}
  \formalstatus{\unformalizedV}
  \bridgestatus{none}
  \notclaimed{This chapter does not claim Riemann integrability, equality of Darboux upper and lower values, host finite-list presentation, quotient partition equality, selected mesh limits, countable choice, ambient RealUp completeness, Lean verification, or bridge checking.}
  \upgradepath{Obligation closure requires checked carrier admission, dyadic-cell coverage, uniform-window compatibility, mesh-refinement exactness, downstream readback nonescape, transport, replay, provenance, and local naming for BEDC.Derived.RiemannPartitionUp.}
\end{closurestatus}
